\section{Coupled coefficients and fixed-basis channel separation}\label{sec:coupled}
This appendix proves Theorem~\ref{thm:separable}. The direct-sum
observation hypothesis is imposed in a specified internal basis; the
coefficient profile deliberately couples those channels in that basis.
The conclusion is not invariant under a change of separability basis.

\subsection{The actual constant profile and its full budget}
Fix $D\ge3$, $n\ge2$, $V=n^D$ and $g=V$. The domain is
$\mathcal U=\mathcal U_0^g$ and its one-channel dimension is
$d_0=(D-1)(V-1)>0$ by \eqref{eq:coexact-dimension}.
There are $h=Dg$ real field components and $r=D^2g$ gradient coordinates.
Let $v=(1,\ldots,1)/\sqrt g\in\R^g$ and put
\begin{equation}\label{eq:coupled-profile}
 \begin{gathered}
 a=c/(16D),\quad H=cn^2,\quad B=H I_g-(H-a)vv^{\mathsf T},\\
 A_{(i,j,\alpha),(i',j',\beta)}
                =\delta_{ii'}\delta_{jj'}B_{\alpha\beta}.
 \end{gathered}
\end{equation}
Here $i$ is the gradient direction, $j$ the spatial one-cochain component,
and $\alpha$ the internal channel. This is a constant positive matrix
with $0<A\le H I$; its actual energy is
\[
 Q_A(z)=\sum_{x,i,j,\alpha,\beta}B_{\alpha\beta}
                D_i^+z_j^\alpha(x)D_i^+z_j^\beta(x).
\]
The weak polarization $v$ has eigenvalue $a$, and $v^\perp$ has eigenvalue
$H$. There are $D^2$ copies of $B$ at each start, so the full local
multiplicities are $D^2$ weak and $D^2(g-1)$ strong. None can be removed
when the tested field domain is restricted.

Since $\sigma/c=\alpha_D^{2/D}\ge9/16$ and $a/c=1/(16D)$, every
weak channel has cap cost one. A strong channel has cost
$\min\{1,\alpha_Dn^{-D}\}$. Thus the exact full budget is
\begin{equation}\label{eq:coupled-budget}
 F_c(A)=D^2V[1+(g-1)\min\{1,\alpha_Dn^{-D}\}]
                    \le D^2(1+\alpha_D)V.
\end{equation}
The minimum is retained in the identity; the inequality uses $g=V=n^D$.
The strong tail is allowed to have order $V$, and is not claimed bounded
independently of volume. Adding $h=DV$ gives $h+F_c(A)\le L_DV$.

\subsection{A projection inside the actual coexact space}
Write the channel-separable observation as
$C=\bigoplus_{\alpha=1}^g C_\alpha$. Each $C_\alpha$ is an arbitrary
finite real linear map on one channel, restricted to $\mathcal U_0$;
within that channel it may mix all sites and spatial components.
Let $R_\alpha$ be the counting orthogonal projection in $\mathcal U_0$
onto the representer space of the restricted rows. This is a projection
within the actual domain, not an unrestricted spatial row-space
projection. Therefore
\[
 R=\bigoplus_\alpha R_\alpha:\mathcal U\to\mathcal U,
 \quad \rank R=r_C:=\rank(C|_{\mathcal U}),\quad
 \ker R=\ker(C|_{\mathcal U}).
\]
The internal direct sum is legitimate. No decomposition of $\mathcal U_0$
by its individual spatial components is assumed.

Define the isometry $T:\mathcal U_0\to\mathcal U$ by
$(Tf)_\alpha=v_\alpha f$. Channel separation gives
\begin{equation}\label{eq:projection-trace}
 T^*RT=g^{-1}\sum_\alpha R_\alpha,\qquad
 \tr(T^*RT)=r_C/g.
\end{equation}
This real symmetric positive contraction has a real unit vector
$f\in\mathcal U_0$ with Rayleigh quotient at most its trace divided by
dimension. As $R$ is an orthogonal projection, there is consequently $f$ with
\[
 t:=\|RTf\|^2\le r_C/(gd_0),\qquad0\le t\le1.
\]
Set $u=(I-R)Tf$. Then $u\in\mathcal U$, $Cu=0$ and
$\|u\|^2=1-t$. The witness may depend on the proposed observation map.
Projections need not preserve locality, gradients, spatial components or
weak polarization; their preservation of $\mathcal U$ is enough.

\subsection{The kernel witness and the necessary rank}
For any periodic field, shifts are isometries, whence
\[
 \sum_i\|D_i^+z\|^2\le4D\|z\|^2.
\]
This counts every original forward start and wrap, including period two,
with no further internal-dimension factor. The constant weak polarization
and the bound $A\le H I$ give, respectively,
\[
 Q_A(Tf)\le4Da,\qquad Q_A(RTf)\le4DHt.
\]
Positivity implies $Q_A(x-y)\le2Q_A(x)+2Q_A(y)$, hence
$Q_A(u)\le8Da+8DHt$. If
\[
 V^{-1}Q_A(z)+\Phi(Cz)\ge c\normn{z}^2\quad(z\in\mathcal U),
 \qquad \Phi(0)=0,
\]
apply this at the actual kernel vector $u$ and multiply by $V$:
\[
 c(1-t)\le8Da+8DHt.
\]
Since $8Da=c/2$ and $H=cn^2$, this yields
\[
 t\ge\frac{c-8Da}{c+8DH}=\frac1{2(1+8Dn^2)}.
\]
Combining with \eqref{eq:projection-trace} proves
\eqref{eq:separable-lower}. No division by $1-t$ occurred; $t=1$ is
harmless. The penalty may be nonlinear, and no differentiability,
quadratic growth or uniform penalty bound was used. Observations remain linear.

For fixed $D\ge3$ the lower-bound expression has order $n^{2D-2}$,
whereas $h+F_c(A)\le L_Dn^D$. Explicitly,
\[
 \frac{r_C}{h+F_c(A)}\ge
 \frac{(D-1)(V-1)}{2L_D(1+8Dn^2)}\longrightarrow\infty.
\]
The divergence has order $n^{D-2}$, including $D=3,4$.
This concerns growing $g=V$, a specified internal basis and finite
unit-spacing tori. It supplies no fixed-$g$ or $D=2$ impossibility and
no physical mass-gap inference.

\subsection{The same profile admits a cheap mixed observation}
Let $Wz=\sum_\alpha v_\alpha z_\alpha$. Since $WT=I$, this map
surjects $\mathcal U$ onto $\mathcal U_0$. Observe its values at every
site and spatial component with rows $V^{-1/2}(Wz)_j(x)$.
The listed output has $DV$ coordinates but actual restricted rank is
exactly $d_0$.
The orthogonal splitting
\[
 z=TWz+z_{\rm strong}
\]
stays in $\mathcal U$, since constant internal combinations commute
with the codifferential. They also commute with forward differences.
The split is orthogonal in norm and in coefficient energy. Every real
component of $z_{\rm strong}$ has zero mean, so the actual torus gap
\eqref{eq:torus-gap} gives
\[
 Q_A(z_{\rm strong})\ge H(16/n^2)\|z_{\rm strong}\|^2
                         =16c\|z_{\rm strong}\|^2.
\]
The weak energy is nonnegative, and
$\|C_{\rm weak}z\|^2=V^{-1}\|Wz\|^2
=V^{-1}\|TWz\|^2$. Therefore
\[
 V^{-1}Q_A(z)+c\|C_{\rm weak}z\|^2\ge c\normn{z}^2.
\]
Its actual rank $d_0=O_D(V)=O_D(h+F_c(A))$ is additive in order on this
family. A fixed global internal rotation exposing $v$ makes the cheap
channel visible. Thus no all-linear or basis-free obstruction follows.
The profile in its original basis is not component-block-diagonal, so
Theorem~\ref{thm:adapted} is not contradicted.
This completes Theorem~\ref{thm:separable}. Appendix~\ref{sec:matching-modes} proves a matching order for this precise family, rather than an all-profile upper bound.
